% =====================================================================
% Closing Group Reflection (250–400 words), placed after Appendices A and B
% =====================================================================
\section*{Closing Group Reflection}
\addcontentsline{toc}{section}{Closing Group Reflection}

AI added the most value in speed and structure. Within minutes it produced stakeholder lists, a checklist of network effects, a complete first canvas and a sensible outline of the cybersecurity problem. Because we ran the prompts in a conversation that already contained our verified research notes, its facts were also largely accurate. Every fee it quoted matched Nord Pool's official fee schedule. This taught us something we had not expected: the accuracy of the output reflected the quality of the input. We cannot conclude from our experience that AI can be trusted with facts it has not been given.

AI failed, or could have misled us, in four recurring ways. First, it presented assumptions as facts, for example the TSOs' preferences and the identity of the ``money side''. It did not distinguish between what was documented and what was inferred. Second, it simplified relationships that the course frameworks show to be more complex. The ``mutual dependency'' between Nord Pool and the TSOs is asymmetric once multi-NEMO competition is taken into account. Third, it slipped into single-firm thinking, most clearly by listing Euronext's futures revenue as Nord Pool's. Fourth, it recommended instruments, such as SLAs and cyber insurance, without checking whether they exist. It also skipped theoretical steps, such as testing whether Nord Pool qualifies as an MSP at all, and it did not link the tasks to each other as the assignment requires.

We also used AI to produce a first draft of the report itself. Revising that draft required the same work as the appendices: checking every claim against a source, asking whose cost, revenue or risk each statement referred to, and rewriting arguments so that they followed from course concepts rather than from generic platform templates.

Our main conclusion is that AI is a fast and useful first drafter and checklist, but not a substitute for economic and strategic judgement. The insights we consider most important were all asymmetries: between contribution and control, between who invests in security and who bears the loss, and between who creates value and who captures it. These insights came from applying course concepts to evidence. We should treat AI output as a hypothesis to be tested, not as an answer.
